\chapter{Code Execution and the Auto-Grading Engine}\label{ch:ch07}

This chapter documents the component at the centre of APEX: the pipeline that takes a
student's source code, runs it safely, judges it fairly, and turns the result into a score.
It is the crown jewel of the system and the piece that justifies the whole project's
existence, because a classroom portal that cannot grade code correctly is just a quiz tool
with a text editor bolted on. Where the rest of APEX is careful CRUD---classes, exams,
enrolments, notifications---the grading engine is where the genuinely hard engineering
lives: running untrusted code without being compromised, comparing free-form program output
without penalising correct answers on cosmetic grounds, and coordinating several
concurrent grading tasks into a single, correct, once-only exam score.

The chapter traces one code submission end to end, from the moment a student presses
\emph{Submit} to the moment the ``Exam Graded'' notification lands, quoting the real
implementation at every stage. It then broadens out to the two other answer formats---
multiple-choice, which is auto-graded by set equality, and written answers, which enter a
manual-grading queue---because a defining claim of APEX is that all three formats flow
through \emph{one} pipeline and aggregate into \emph{one} attempt score. The whole flow is
summarised in \cref{fig:grading-seq}; the sections below expand each box of that figure into
its underlying code.

Two design principles recur throughout and are worth stating up front. First,
\textbf{APEX never runs untrusted code on its own host}: every execution is delegated over
HTTP to Judge0~\cite{judge0,dosilovic2018judge0}, which runs the code inside the
\texttt{isolate} sandbox~\cite{isolate}. Second, \textbf{APEX judges output by content, not
by bytes}: a nine-line normalisation routine strips away the whitespace differences that
cause the most common false negatives in output-comparison grading, a fairness concern the
auto-assessment literature has flagged for two decades~\cite{douce2005autoassess}. These two
ideas---delegate execution, own the judging---shape everything that follows.

\section{Execution Model: Delegating to a Sandbox}\label{sec:grading-sandbox}

\subsection{Why the sandbox is not ours}

Running a student's arbitrary program is the single most dangerous thing an assessment
platform does. A submission is, by definition, untrusted code: it may loop forever, allocate
until the machine swaps to death, fork-bomb the process table, try to read the server's
secrets, or open a socket to exfiltrate data. Executing such code directly on the
application host---whether through \texttt{os/exec} or a hand-rolled container---makes the
grader itself the attack surface, and writing a correct sandbox from scratch is one of the
hardest problems in systems security. APEX declines that problem entirely.

Instead, execution is delegated to \textbf{Judge0}~\cite{judge0}, an open-source
``code execution as a service'' whose design and rationale are described by da Silva Dos Reis and others~\cite{dosilovic2018judge0}.
Judge0 accepts a JSON document of \texttt{\{language, source\_code, stdin\}} over HTTP,
compiles and runs the code, and returns \texttt{\{stdout, stderr, compile\_output, status,
time, memory\}}. Crucially, Judge0 runs each submission inside
\textbf{\texttt{isolate}}~\cite{isolate}, the sandbox originally built for the International
Olympiad in Informatics to grade contestants' programs safely. \texttt{isolate} uses Linux
namespaces and control groups (\texttt{cgroups}) to give the untrusted process an isolated
view of the filesystem and network and to cap its CPU time and memory. The security posture
follows directly: the APEX backend sits \emph{outside} the sandbox and communicates with it
only by an HTTP request and response. Untrusted student code never touches the application
host, and the most security-critical component in the system is the one piece the team
deliberately did not write.

The trade-off is deliberate and accepted: grading now depends on a network hop and on Judge0's
availability. If Judge0 is unreachable, execution fails and the caller receives an explicit
error rather than a wrong grade (\Cref{sec:grading-client}). The team judged that
correctness and safety outweigh self-sufficiency---a defensible position for a grading
system, where a wrong ``accepted'' verdict is far worse than a temporary outage.

By default the backend targets the public endpoint \texttt{https://ce.judge0.com}, but the
base URL is configurable through the \texttt{JUDGE0\_URL} environment variable and wired at
start-up via \texttt{judge0.Init(cfg.Judge0URL)}. The public Community Edition endpoint is
rate-limited and not intended for production; a real deployment self-hosts Judge0. Note that
Judge0 is reached as an \emph{external} HTTP service and is therefore not one of the four
Docker Compose services described in \cref{ch:ch10}.

\subsection{The language map}\label{sec:grading-langs}

APEX supports six execution languages. The mapping from the language string sent by the
frontend to Judge0's numeric language identifier is a single map in
\texttt{internal/judge0/client.go}, reproduced in \cref{lst:langmap}. Several human-friendly
aliases collapse onto the same identifier---\texttt{python} and \texttt{python3} both resolve
to 100, and \texttt{cpp} and \texttt{c++} both resolve to 105---so the six distinct
languages are exposed through ten map keys.

\begin{lstlisting}[language=Go,caption={The language-to-Judge0-ID map (\texttt{client.go}).},label={lst:langmap}]
var LanguageMap = map[string]int{
	"python":     100,
	"python3":    100,
	"javascript": 102,
	"js":         102,
	"typescript": 101,
	"ts":         101,
	"java":       91,
	"c":          103,
	"cpp":        105,
	"c++":        105,
}
\end{lstlisting}

The six languages and their identifiers are listed in \cref{tab:grading-langs}. A language
string absent from the map is not silently passed to Judge0: \texttt{RunCode} returns an
internal sentinel result (status identifier \texttt{-1}) carrying the message
``\texttt{unsupported language}'', and the HTTP handlers additionally pre-validate the
language against \texttt{LanguageMap} before any network call is made, so an unknown
language is rejected with a \texttt{400 Bad Request} rather than a gateway error.

\begin{table}[htbp]
\centering
\caption{Supported execution languages and their Judge0 language identifiers.}
\label{tab:grading-langs}
\begin{tabular}{@{}L{4.5cm}L{4cm}c@{}}
\toprule
\textbf{Language} & \textbf{Accepted keys} & \textbf{Judge0 ID} \\
\midrule
Python     & \texttt{python}, \texttt{python3} & 100 \\
JavaScript & \texttt{javascript}, \texttt{js}   & 102 \\
TypeScript & \texttt{typescript}, \texttt{ts}   & 101 \\
Java       & \texttt{java}                      & 91  \\
C          & \texttt{c}                         & 103 \\
C\texttt{++} & \texttt{cpp}, \texttt{c++}       & 105 \\
\bottomrule
\end{tabular}
\end{table}

At the user-interface layer the coverage is deliberately asymmetric, which is worth
stating precisely. The standalone live-coding playground exposes all six languages for
free experimentation, whereas the exam-builder's coding-question editor currently offers
four of them---Python, JavaScript, C, and C\texttt{++}---as authoring targets. This is a
front-end configuration choice, not a limitation of the grader: the execution and grading
backend accepts any of the six keys in \cref{tab:grading-langs} regardless of which surface
produced the submission, so widening the exam-builder set is a one-line change to the editor
rather than any change to the pipeline described in this chapter.

\subsection{The Judge0 client}\label{sec:grading-client}

Every execution in APEX---playground runs, sample-test feedback, and real grading---funnels
through one function, \texttt{RunCode} in \texttt{internal/judge0/client.go}. It marshals the
request, performs a single HTTP \texttt{POST}, and decodes the response into a compact
\texttt{CodeResult}. The core is shown in \cref{lst:runcode}.

\begin{lstlisting}[language=Go,caption={The single point of contact with the sandbox (\texttt{client.go}).},label={lst:runcode}]
func RunCode(code, language, stdin string) CodeResult {
	langID, ok := LanguageMap[language]
	if !ok {
		return CodeResult{StatusID: -1, Stderr: "unsupported language: " + language}
	}

	reqBody := request{LanguageID: langID, SourceCode: code, Stdin: stdin}
	body, err := json.Marshal(reqBody)
	if err != nil {
		return CodeResult{StatusID: -1, Stderr: "failed to marshal request"}
	}

	url := BaseURL + "/submissions?wait=true&base64_encoded=false&fields=stdout,stderr,compile_output,status,time,memory"
	client := &http.Client{Timeout: 30 * time.Second}
	resp, err := client.Post(url, "application/json", bytes.NewReader(body))
	if err != nil {
		return CodeResult{StatusID: -1, Stderr: "judge0 unreachable"}
	}
	defer resp.Body.Close()
	// ... read body, unmarshal, parse time (seconds -> ms) and memory (kb) ...
}
\end{lstlisting}

Several details of the request are deliberate and defensible:

\begin{itemize}[leftmargin=1.4em]
  \item \textbf{Synchronous execution (\texttt{wait=true}).} The single \texttt{POST} blocks
  until Judge0 has finished compiling and running the code, and the verdict comes back in the
  same response. There is no client-side poll loop and no callback URL to manage. A
  \textbf{30-second} client timeout guards against a hung sandbox; if it fires, \texttt{RunCode}
  returns the \texttt{-1} sentinel.
  \item \textbf{Plain text, not base64 (\texttt{base64\_encoded=false}).} Source code and
  standard input are sent as raw JSON strings. Judge0 supports base64 transport, but APEX opts
  out for simplicity; the modest cost is that exotic byte sequences could in principle be
  mangled, which is acceptable for classroom source code.
  \item \textbf{Field projection (\texttt{fields=...}).} The response is trimmed to exactly the
  six fields APEX consumes, keeping the payload small.
  \item \textbf{Metric extraction.} Judge0 reports execution time as a seconds string and
  memory as a number; \texttt{RunCode} parses the former into integer milliseconds
  (\texttt{seconds}$\times 1000$) and the latter into kilobytes, populating
  \texttt{CodeResult.TimeMs} and \texttt{CodeResult.MemoryKb}. These are \emph{measured} and
  recorded, a point revisited in \Cref{sec:grading-limits}.
\end{itemize}

The status identifiers Judge0 returns drive the rest of the pipeline. \Cref{tab:grading-status}
lists the ones APEX interprets. Two entries deserve emphasis now and are explained fully in
\Cref{sec:grading-loop}: status \texttt{3} (Accepted) and status \texttt{4} (Wrong Answer)
are \emph{both} treated by APEX as ``the program ran cleanly,'' because APEX performs its own
output comparison and does not defer to Judge0's built-in diff.

\begin{table}[htbp]
\centering
\caption{Judge0 status identifiers and APEX's interpretation.}
\label{tab:grading-status}
\begin{tabular}{@{}cL{5.2cm}L{4.6cm}@{}}
\toprule
\textbf{ID} & \textbf{Judge0 meaning} & \textbf{APEX action} \\
\midrule
3  & Accepted (ran cleanly)      & Counts as ``ran'' $\rightarrow$ compare normalised output \\
4  & Wrong Answer (Judge0 diff)  & \emph{Also} counts as ``ran'' $\rightarrow$ re-compare with APEX's own normalised diff \\
5  & Time Limit Exceeded         & Verdict \texttt{time\_limit\_exceeded} \\
6  & Compilation Error           & Verdict \texttt{compilation\_error} \\
11 & Runtime Error               & Verdict \texttt{runtime\_error} \\
-1 & APEX sentinel (unreachable / marshal / parse failure) & Playground returns \texttt{504}; other paths fail the test \\
\bottomrule
\end{tabular}
\end{table}

\section{Three Entry Points, One That Grades}\label{sec:grading-entrypoints}

A student's code can reach the backend through three distinct endpoints, and keeping them
separate is one of the cleaner decisions in the system. Only one of the three persists a
graded \texttt{Submission}; the other two are ephemeral feedback paths that write nothing to
the database. \Cref{tab:grading-entrypoints} summarises the split.

\begin{table}[htbp]
\centering
\caption{The three code-execution entry points.}
\label{tab:grading-entrypoints}
\begin{tabular}{@{}L{4.6cm}L{4.4cm}cL{3.4cm}@{}}
\toprule
\textbf{Endpoint} & \textbf{Handler} & \textbf{Persists grade?} & \textbf{Purpose} \\
\midrule
\texttt{POST /execute}                & \texttt{execute.Execute}      & No  & Playground scratchpad \\
\texttt{POST /submissions/run}        & \texttt{submission.RunSolution} & No & Sample-test feedback \\
\texttt{POST /student/exams/:id/submit} & \texttt{exam.SubmitAttempt}   & Yes & The real, graded submission \\
\bottomrule
\end{tabular}
\end{table}

The practical consequence is that a student can press ``Run'' fifty times during an exam at a
cost of zero database rows and zero grading work. Running code is feedback; submitting is the
grade; they are different code paths on purpose.

\subsection{The playground path}

The cheapest path is the playground handler, \texttt{execute.Execute} in
\texttt{internal/execute/handler.go}. It binds a \texttt{\{language, code, stdin\}} request,
pre-validates the language against \texttt{LanguageMap}, calls \texttt{RunCode}, and returns
the raw output. The one subtlety is its error mapping: the internal \texttt{-1} sentinel is
surfaced as an HTTP \texttt{504 Gateway Timeout}, a precise signal that \emph{the sandbox
service} failed rather than that \emph{the student's code} was wrong.

\begin{lstlisting}[language=Go,caption={The playground handler maps the sandbox sentinel to 504 (\texttt{execute/handler.go}).},label={lst:execute}]
if _, ok := judge0.LanguageMap[req.Language]; !ok {
	c.JSON(http.StatusBadRequest, gin.H{"error": "unsupported language: " + req.Language})
	return
}

result := judge0.RunCode(req.Code, req.Language, req.Stdin)
if result.StatusID == -1 {
	c.JSON(http.StatusGatewayTimeout, gin.H{"error": result.Stderr})
	return
}
c.JSON(http.StatusOK, gin.H{
	"stdout": result.Stdout, "stderr": result.Stderr,
	"status_id": result.StatusID, "time_ms": result.TimeMs, "memory_kb": result.MemoryKb,
})
\end{lstlisting}

\subsection{The sample-test feedback path}

The second non-persisting path is \texttt{submission.RunSolution}
(\texttt{internal/submission/handler.go}), which powers the ``Run sample tests'' button during
an exam. It first enforces resource ownership---the third authorisation gate described in
\cref{ch:ch06}: a teacher may run only against a problem they own, and a student may run only
against a problem belonging to an exam assigned to a class they are enrolled in, verified by a
join across \texttt{exam\_classes} and \texttt{class\_members}. It then loads only the
\emph{sample} test cases (\texttt{WHERE is\_sample = true}), ordered by \texttt{order\_index},
and evaluates them through \texttt{RunAgainstTestCases}. Nothing is written; the student
receives per-test input, expected output, and actual output so they can debug before
committing to a real submission. Because \texttt{RunAgainstTestCases} applies the \emph{same}
\texttt{normalizeOutput} routine as the real grader (\Cref{sec:grading-normalize}), the
feedback a student sees while iterating matches the verdict they will later be graded by.

\section{The Coding Grader}\label{sec:grading-coding}

The graded path begins in \texttt{exam.SubmitAttempt}, which creates one \texttt{Submission}
row per answer at status \texttt{pending} and dispatches a grader according to the problem
type. For a coding problem, that grader is \texttt{judge0.Grade}. The full sequence---from the
Monaco editor through this grader to the final notification---is depicted in
\cref{fig:grading-seq}.

\begin{figure}[htbp]
\centering
\resizebox{\textwidth}{!}{%
\begin{tikzpicture}[
  font=\sffamily\scriptsize,
  head/.style={draw=apexblue, fill=apexlight, rounded corners=2pt,
    align=center, font=\sffamily\scriptsize\bfseries, minimum height=8mm,
    inner sep=3pt, text width=20mm},
  headext/.style={head, draw=apexamber, fill=apexamber!12},
  headstore/.style={head, draw=apexgreen!70!black, fill=apexgreen!8},
  life/.style={apexgrey, dashed, thick},
  call/.style={-{Latex[length=2mm]}, draw=apexblue, thick},
  callg/.style={-{Latex[length=2mm]}, draw=apexgrey, thick},
  ret/.style={-{Latex[length=2mm]}, draw=apexgrey, densely dashed},
  msg/.style={apexlabel, font=\sffamily\scriptsize},
]

% --- lifeline x-coordinates (evenly spaced) --------------------------
\def\xA{0}    % Student SPA
\def\xB{3.6}  % Gin Handler
\def\xC{7.2}  % judge0.Grade goroutine
\def\xD{10.8} % Judge0 CE
\def\xE{14.0} % isolate sandbox
\def\xF{17.2} % PostgreSQL

\def\ytop{0}
\def\ybot{-11.4}

% --- head boxes ------------------------------------------------------
\node[head]      (hA) at (\xA,\ytop) {Student\\SPA};
\node[head]      (hB) at (\xB,\ytop) {Gin\\Handler};
\node[head]      (hC) at (\xC,\ytop) {\texttt{judge0.Grade}\\goroutine};
\node[headext]   (hD) at (\xD,\ytop) {Judge0 CE};
\node[headext]   (hE) at (\xE,\ytop) {\texttt{isolate}\\sandbox};
\node[headstore] (hF) at (\xF,\ytop) {PostgreSQL};

% --- lifelines -------------------------------------------------------
\foreach \x in {\xA,\xB,\xC,\xD,\xE,\xF}
  \draw[life] (\x,-0.5) -- (\x,\ybot);

% --- message rows (decreasing y) -------------------------------------
\def\ya{-1.3}  % SPA -> Handler
\def\yb{-2.2}  % Handler -> DB pending
\def\yc{-3.1}  % Handler -> goroutine (spawn)
\def\yd{-4.0}  % Handler -> SPA (202)
% loop body
\def\ye{-5.5}  % goroutine -> Judge0
\def\yf{-6.4}  % Judge0 -> isolate
\def\yg{-7.3}  % isolate -> Judge0 (return)
\def\yh{-8.2}  % Judge0 -> goroutine (return)
\def\yi{-9.1}  % goroutine self: normalize + diff
% post loop
\def\yj{-10.2} % goroutine -> DB (TestResults + Submission)
\def\yk{-11.0} % goroutine -> DB (ExamAttempt)

% 1. submit request
\draw[call] (\xA,\ya) -- (\xB,\ya)
  node[msg,midway,above]{\texttt{POST /student/exams/:id/submit}};

% 2. persist pending submission
\draw[call] (\xB,\yb) -- (\xF,\yb)
  node[msg,midway,above]{INSERT Submission (\texttt{status=pending})};

% 3. spawn grading goroutine (fire-and-forget)
\draw[callg] (\xB,\yc) -- (\xC,\yc)
  node[msg,midway,above]{\texttt{go judge0.Grade(...)}\;\emph{(fire-and-forget)}};

% 4. immediate response
\draw[ret] (\xB,\yd) -- (\xA,\yd)
  node[msg,midway,above]{\textasciitilde202 Accepted (returns at once)};

% --- loop frame over the test cases ----------------------------------
\begin{scope}[on background layer]
  \fill[apexblue!4,rounded corners=2pt] (\xC-1.0,-4.9) rectangle (\xF+0.9,-9.55);
\end{scope}
\draw[apexblue!55,rounded corners=2pt] (\xC-1.0,-4.9) rectangle (\xF+0.9,-9.55);
\node[fill=apexblue!55,text=white,font=\sffamily\scriptsize\bfseries,
      anchor=north west,inner sep=2pt] at (\xC-1.0,-4.9)
  {loop\; [for each \texttt{TestCase} by \texttt{order\_index}]};

% 5. send code+stdin to Judge0
\draw[call] (\xC,\ye) -- (\xD,\ye)
  node[msg,midway,above]{\texttt{POST /submissions?wait=true}\;(code, lang, stdin)};

% 6. run inside isolate
\draw[call] (\xD,\yf) -- (\xE,\yf)
  node[msg,midway,above]{run under \texttt{time\_limit\_ms} / \texttt{memory\_limit\_kb}};

% 7. isolate returns
\draw[ret] (\xE,\yg) -- (\xD,\yg)
  node[msg,midway,above]{stdout + exit status};

% 8. Judge0 returns
\draw[ret] (\xD,\yh) -- (\xC,\yh)
  node[msg,midway,above]{stdout + status id};

% 9. self message: normalize + diff
\draw[callg] (\xC,\yi-0.18) -- ++(1.35,0) -- ++(0,0.36) -- (\xC,\yi+0.18);
\node[msg,anchor=west] at (\xC+0.15,\yi+0.55)
  {\texttt{normalizeOutput} + diff vs expected \(\rightarrow\) pass/fail, write \texttt{TestResult}};

% 10. aggregate score, update submission
\draw[call] (\xC,\yj) -- (\xF,\yj)
  node[msg,midway,above]{UPDATE Submission: \texttt{score = passed/total*100}, status};

% 11. roll into attempt
\draw[call] (\xC,\yk) -- (\xF,\yk)
  node[msg,midway,above]{\texttt{maybeFinalizeAttempt}: roll score into ExamAttempt};

\end{tikzpicture}%
}
\caption{Sequence of an auto-graded code submission.}
\label{fig:grading-seq}
\end{figure}


\subsection{The fire-and-forget dispatch}\label{sec:grading-dispatch}

Each submission is graded in its own bare goroutine. \texttt{SubmitAttempt} creates the row and
then launches the appropriate grader with the \texttt{go} keyword, routing by problem type:

\begin{lstlisting}[language=Go,caption={Per-submission grader dispatch in \texttt{SubmitAttempt}.},label={lst:dispatch}]
switch submissionType {
case "mcq":
	go judge0.GradeMCQ(sub.ID)
case "written":
	go judge0.GradeWritten(sub.ID)
default:
	go judge0.Grade(sub.ID)
}
\end{lstlisting}

The \texttt{go} keyword means the HTTP handler does not wait for grading to finish: the request
returns \texttt{201 Created} immediately, carrying the attempt identifier and the list of
created submission identifiers, while grading proceeds in the background. The student sees an
``Exam Submitted'' notification at once; the ``Exam Graded'' notification arrives later, once
every goroutine has completed.

This model is deliberately simple and it works at the concurrency of a single classroom, but
it carries a durability gap that the project documents openly (\texttt{README.md} line
39). The dispatch is \textbf{fire-and-forget}: there is no work queue, no supervisor, and no
retry. As the next subsection shows, \texttt{Grade} immediately flips the submission to
\texttt{running}; if the backend process crashes or is redeployed mid-grade, that goroutine
dies with the process and the row is left stuck in \texttt{running} forever. Because the
attempt finaliser refuses to proceed while any sibling submission is still
\texttt{pending} or \texttt{running} (\Cref{sec:grading-aggregate}), the attempt never
finalises and the student never receives the ``Exam Graded'' notification. The failure mode is
discussed in full, together with its proposed fix, in \Cref{sec:grading-queue}.

\subsection{Loading test cases and the empty-set edge case}

\texttt{Grade} begins by loading the submission, flipping its status to \texttt{running}, and
loading \emph{all} test cases for the problem ordered by \texttt{order\_index}. There is
deliberately no \texttt{is\_sample} filter here---this is the real grade, so both sample and
hidden tests run. One edge case is handled defensively: a problem with \emph{no} test cases
configured is auto-accepted at score \texttt{100.0}, on the principle that a teacher who forgot
to add tests should not thereby fail the student.

\begin{lstlisting}[language=Go,caption={Grader set-up and the no-test-case edge case (\texttt{grader.go}).},label={lst:grade-setup}]
database.DB.Model(&submission).Update("status", "running")

var testCases []models.TestCase
database.DB.Where("problem_id = ?", submission.ProblemID).Order("order_index asc").Find(&testCases)

if len(testCases) == 0 {
	database.DB.Model(&submission).Updates(map[string]any{
		"status": "accepted",
		"score":  100.0,
	})
	return
}
\end{lstlisting}

\subsection{The grading loop}\label{sec:grading-loop}

The heart of the coding grader is the loop that runs the student's code against each test case
and records a per-test verdict. It is quoted in full in \cref{lst:grade-loop}.

\begin{lstlisting}[language=Go,caption={The per-test-case grading loop (\texttt{grader.go}).},label={lst:grade-loop}]
for _, tc := range testCases {
	result := RunCode(submission.Code, submission.Language, tc.Input)

	actualOutput := normalizeOutput(result.Stdout)
	expectedOutput := normalizeOutput(tc.ExpectedOutput)
	// StatusID 3 = Accepted; treat StatusID 4 (Wrong Answer) as valid-run-but-compare.
	ran := result.StatusID == 3 || result.StatusID == 4
	isPassed := ran && actualOutput == expectedOutput

	if isPassed {
		passed++
	}

	status := "accepted"
	if !isPassed {
		switch result.StatusID {
		case 6:
			status = "compilation_error"
		case 5:
			status = "time_limit_exceeded"
		case 11:
			status = "runtime_error"
		default:
			status = "wrong_answer"
		}
	}

	tr := models.TestResult{
		SubmissionID:    submission.ID,
		TestCaseID:      tc.ID,
		Passed:          isPassed,
		ActualOutput:    actualOutput,
		ExecutionTimeMs: result.TimeMs,
		MemoryKb:        result.MemoryKb,
		Status:          status,
	}
	database.DB.Create(&tr)

	if result.TimeMs > maxTimeMs {
		maxTimeMs = result.TimeMs
	}
	if result.MemoryKb > maxMemoryKb {
		maxMemoryKb = result.MemoryKb
	}
}
\end{lstlisting}

For each test case, the loop performs four conceptual steps. First, it sends the student's code
and this test's input to Judge0 through \texttt{RunCode}. Second, it passes \emph{both} the
program's actual standard output \emph{and} the test's expected output through the same
\texttt{normalizeOutput} routine---applied symmetrically, which is the fairness guarantee
analysed in \Cref{sec:grading-normalize}. Third, it decides whether the test passed. Fourth,
it writes one \texttt{TestResult} row recording the outcome, the actual output, and the
measured time and memory, and it tracks the maximum time and memory seen across all tests.

The pass decision hides a subtle and important design choice. The predicate
\texttt{ran := result.StatusID == 3 || result.StatusID == 4} treats \emph{both} Judge0's
``Accepted'' (3) and its ``Wrong Answer'' (4) as ``the program executed cleanly.'' The reason
is that Judge0 performs its \emph{own} byte-exact diff and may return status 4 purely because
of a trailing newline---exactly the false negative APEX refuses to honour. So APEX ignores
Judge0's accept/reject decision and re-judges with its own normalised equality check:
\texttt{isPassed := ran \&\& actualOutput == expectedOutput}. In effect, \textbf{Judge0 is
treated as an executor---run this safely and tell me what it printed---not as a judge}. The
judging belongs to APEX. Only genuine execution failures (compilation error 6, time-limit 5,
runtime error 11) are surfaced as failure verdicts; anything else that ran but did not match is
recorded as \texttt{wrong\_answer}.

\subsection{Output normalisation: nine lines of fairness}\label{sec:grading-normalize}

Output-comparison grading has one notorious weakness, identified in the literature as far back
as the first systematic surveys of automated assessment~\cite{douce2005autoassess}: a program
that computes the \emph{right} answer can be marked wrong because of invisible whitespace. A
student prints \texttt{42\textbackslash n} where the expected answer is \texttt{42} with no
trailing newline; a student on Windows emits \texttt{\textbackslash r\textbackslash n} line
endings where the teacher typed \texttt{\textbackslash n}; a stray trailing space follows a
correct number. Under naive byte-for-byte comparison every one of these is a false negative.
This is the single most common unfairness on any output-diffing platform, and eliminating it is
the entire job of \texttt{normalizeOutput}, reproduced in full in \cref{lst:normalize}.

\begin{lstlisting}[language=Go,caption={\texttt{normalizeOutput}---applied symmetrically to actual and expected output (\texttt{grader.go}).},label={lst:normalize}]
// normalizeOutput strips CR, trailing whitespace on each line, and trims edges
// so Judge0 output (often \r\n with trailing newline) compares cleanly to
// teacher-entered expected output.
func normalizeOutput(s string) string {
	s = strings.ReplaceAll(s, "\r\n", "\n")
	s = strings.ReplaceAll(s, "\r", "\n")
	lines := strings.Split(s, "\n")
	for i, ln := range lines {
		lines[i] = strings.TrimRight(ln, " \t")
	}
	return strings.TrimSpace(strings.Join(lines, "\n"))
}
\end{lstlisting}

The routine applies four normalisation rules, in order, and each one closes a specific class of
false negative:

\begin{enumerate}[leftmargin=1.6em]
  \item \texttt{ReplaceAll("\textbackslash r\textbackslash n", "\textbackslash n")} converts
  Windows CRLF line endings to Unix LF, so a student typing on Windows---or a teacher who pasted
  expected output out of a Windows editor---does not fail on invisible carriage returns.
  \item \texttt{ReplaceAll("\textbackslash r", "\textbackslash n")} folds any lone carriage
  return (classic-Mac line endings) into a newline as well: belt and suspenders.
  \item Splitting on \texttt{\textbackslash n} and applying
  \texttt{TrimRight(ln, " \textbackslash t")} to \emph{each} line strips trailing spaces and
  tabs from every line, so a stray space after a printed value is not treated as a correctness
  error.
  \item \texttt{TrimSpace} on the whole joined string removes leading and trailing whitespace
  from the entire output, absorbing the final trailing newline that \texttt{print} and
  \texttt{println} routinely append, as well as any leading blank line.
\end{enumerate}

The fairness guarantee rests on one property: because the \emph{same} function is applied to
both the actual output and the expected output before comparison, the comparison is symmetric.
It measures the \emph{correctness of content}, not the \emph{byte-exactness of whitespace}.
Without these nine lines, a student who computed the right answer would routinely fail on a
trailing newline, a CRLF line ending, or a trailing space---and each of those is a formatting
artefact, not a wrong answer.

It is equally important to state plainly what this normalisation does \emph{not} do. It is
whitespace-trim-fair, not token-fair: it does not collapse \emph{internal} runs of spaces, so
\texttt{1  2} and \texttt{1 2} still differ, and it carries no floating-point tolerance, so
\texttt{3.14} and \texttt{3.140000001} still differ. It solves the most common false negative,
not every conceivable one. That is a deliberate scoping decision rather than an oversight: the
routine is small enough to verify by eye, and its behaviour is predictable, which for a
fairness-critical function is itself a virtue.

\subsection{Scoring and aggregation into the submission}\label{sec:grading-score}

Once every test case has been run, the grader computes a proportional score and writes the
outcome back to the \texttt{Submission} row:

\begin{lstlisting}[language=Go,caption={Coding score and final submission update (\texttt{grader.go}).},label={lst:grade-score}]
score := float64(passed) / float64(total) * 100.0
finalStatus := "wrong_answer"
if passed == total {
	finalStatus = "accepted"
}

database.DB.Model(&submission).Updates(map[string]any{
	"status":            finalStatus,
	"passed_count":      passed,
	"total_count":       total,
	"score":             score,
	"execution_time_ms": maxTimeMs,
	"memory_kb":         maxMemoryKb,
})
\end{lstlisting}

Coding submissions therefore earn \textbf{partial credit}: passing seven of ten tests yields a
score of \texttt{70.0}, and the submission's \texttt{status} is \texttt{accepted} only if every
test passed, otherwise \texttt{wrong\_answer}. The counts \texttt{passed\_count} and
\texttt{total\_count} are stored for display, and the \emph{worst-case} (maximum) execution time
and memory across all tests are recorded on the submission. Notably, the per-test-case
\texttt{points} field that exists in the \texttt{TestCase} schema is \emph{not} used by the
coding grader---coding is scored as a flat \texttt{passed / total} fraction, and points
weighting happens one level up, at the problem level, during attempt aggregation
(\Cref{sec:grading-aggregate}).

Finally, if the submission is standalone---not part of an exam attempt, i.e.
\texttt{ExamAttemptID == nil}, as with a bank-question practice run---a per-submission ``result''
notification fires immediately. If instead the submission belongs to an exam attempt, that
per-submission notification is suppressed and control passes to \texttt{maybeFinalizeAttempt}, so
the student receives \emph{one} exam-level notification rather than one per question.

\subsection{Limits: measured but not yet pushed}\label{sec:grading-limits}

The \texttt{Problem} model carries two resource-limit fields---\texttt{time\_limit\_ms}
(default \texttt{2000}) and \texttt{memory\_limit\_kb} (default \texttt{262144}, i.e.\ 256~MB)---
intended to let a teacher tighten the sandbox per problem. In the current implementation these
fields are \emph{not} forwarded to Judge0: \texttt{RunCode} sends no \texttt{cpu\_time\_limit}
or \texttt{memory\_limit}, so the effective limits are Judge0's own defaults. APEX
\emph{measures} time and memory from each response and records them on every \texttt{TestResult}
and \texttt{Submission}, but it does not yet \emph{push} the per-problem ceilings down. The
system remains safe---a runaway loop still becomes a time-limit-exceeded verdict under Judge0's
default ceiling---but a teacher who sets a one-second limit does not presently get it enforced.
Wiring these two fields into the request body is a small, known task rather than a design flaw.

\section{The Other Two Formats}\label{sec:grading-other}

APEX's defining structural claim is that coding, multiple-choice, and written answers all flow
through one pipeline and aggregate into one attempt score. The two non-coding graders live
alongside \texttt{Grade} in \texttt{grader.go} and share its dispatch, notification, and
finalisation machinery.

\subsection{Multiple-choice: auto-graded by set equality}\label{sec:grading-mcq}

\texttt{GradeMCQ} judges a multiple-choice answer by \textbf{exact set equality} between the
teacher's correct option identifiers and the student's selected option identifiers. Both are
stored as JSON arrays---\texttt{Problem.CorrectOptionIDs} and
\texttt{Submission.SelectedOptions}---which the grader unmarshals into two Go sets, trimming and
dropping empty entries. The answer is correct only if the two sets are exactly equal:
same cardinality, and every correct identifier selected.

\begin{lstlisting}[language=Go,caption={Set-equality decision in \texttt{GradeMCQ} (\texttt{grader.go}).},label={lst:mcq}]
isCorrect := len(correctSet) == len(selectedSet)
if isCorrect {
	for id := range correctSet {
		if !selectedSet[id] {
			isCorrect = false
			break
		}
	}
}
\end{lstlisting}

MCQ grading is \textbf{all-or-nothing}: a correct answer scores \texttt{100.0} with status
\texttt{accepted}, an incorrect one scores \texttt{0.0} with status \texttt{wrong\_answer}, and
there is no partial credit for selecting a subset of a multi-answer question. One protective
special case precedes the comparison: if a teacher has configured \emph{no} correct answers
(\texttt{len(correctSet) == 0}), the submission is marked \texttt{pending\_review} at score
\texttt{0} rather than failed, on the principle that a student must not be failed on a question
whose answer key was never set. A JSON parse failure on either field marks the submission
\texttt{error}. As with coding, standalone submissions fire an immediate notification, while
exam submissions defer to \texttt{maybeFinalizeAttempt}.

\subsection{Written answers: the manual-grading queue}\label{sec:grading-written}

\texttt{GradeWritten} performs no automatic judgement at all. It records the student's answer,
sets the submission to \texttt{pending\_review} at score \texttt{0} with \texttt{total\_count} of
1, and calls \texttt{maybeFinalizeAttempt}. In effect it enrols the answer into a
\emph{manual-grading queue}: a written answer never auto-fails a student; it always waits for a
human. This is the correct default for free-form prose, which no output diff can grade.

Teachers drain that queue through two endpoints in \texttt{internal/submission/handler.go}.
\texttt{GetSubmission} lets an authorised reader---the student who owns the submission, or the
teacher who owns the exam---retrieve a submission with its per-test results; notably, it exposes
the \texttt{input} and \texttt{expected\_output} of a test \emph{only} when that test is a
sample (\texttt{if rv.IsSample}), so hidden test cases are never leaked to students even on the
results screen. \texttt{GradeSubmission} handles the actual manual grade via
\texttt{PUT /submissions/:id/grade}, gated to the exam owner:

\begin{lstlisting}[language=Go,caption={Manual grading applies a score override and re-finalises the attempt (\texttt{submission/handler.go}).},label={lst:manualgrade}]
if exam.TeacherID != teacherID {
	c.JSON(http.StatusForbidden, gin.H{"error": "not authorized to grade this submission"})
	return
}

sub.Score = req.Score
sub.Status = req.Status
sub.PassedCount = int(req.Score)
sub.TotalCount = 100
if req.TeacherFeedback != nil {
	sub.TeacherFeedback = *req.TeacherFeedback
}
database.DB.Save(&sub)

if sub.ExamAttemptID != nil {
	judge0.FinalizeAttempt(*sub.ExamAttemptID)
}
\end{lstlisting}

The teacher supplies a \texttt{score}, a \texttt{status}, and optional
\texttt{teacher\_feedback}, and the handler writes them onto the submission. Crucially it then
calls \texttt{judge0.FinalizeAttempt}, the public wrapper over \texttt{maybeFinalizeAttempt}, so
that grading a written answer immediately recomputes the whole attempt score and---if it was the
last outstanding answer---fires the ``Exam Graded'' notification. This is exactly why the
aggregator is written as a re-summation rather than an incremental update: a manual grade lands
long after the automatic ones, and re-aggregating from the children makes that late arrival
correct without any special-casing.

\section{Attempt Aggregation}\label{sec:grading-aggregate}

Every grader---coding, MCQ, and written---ends by calling \texttt{maybeFinalizeAttempt}. This
function is the orchestrator that turns many per-question submissions into one exam score and one
notification, coordinating several concurrent goroutines without a lock.

\subsection{The barrier: only the last goroutine proceeds}

The function's first act is a guard. It counts sibling submissions still in \texttt{pending} or
\texttt{running} state and returns immediately if any remain:

\begin{lstlisting}[language=Go,caption={The concurrency barrier in \texttt{maybeFinalizeAttempt} (\texttt{grader.go}).},label={lst:barrier}]
var pending int64
database.DB.Model(&models.Submission{}).
	Where("exam_attempt_id = ? AND status IN ?", *attemptID, []string{"pending", "running"}).
	Count(&pending)
if pending > 0 {
	return
}
\end{lstlisting}

Several goroutines, each grading one answer, all call this same function as they finish, but only
the \emph{last} one to complete finds zero pending submissions and proceeds past the guard. This
is how independent goroutines coordinate into a single finalisation without a mutex or a shared
counter: the database row-count \emph{is} the coordination primitive. It also explains the
finalisation half of the fire-and-forget gap---a submission stuck in \texttt{running} because its
goroutine died keeps this count above zero forever, and the attempt never finalises.

\subsection{Re-summing over the exam's full problem list}

Past the barrier, the aggregator recomputes the attempt score from scratch. It loads the
\emph{exam's full problem list}---not the submissions---and folds each problem's contribution
into an earned/total ratio:

\begin{lstlisting}[language=Go,caption={Points-weighted re-aggregation over the problem list (\texttt{grader.go}).},label={lst:aggregate}]
for _, p := range problems {
	s, hasSub := subByProblem[p.ID]
	if hasSub && s.Status == "pending_review" {
		continue
	}
	pts := float64(p.Points)
	if pts <= 0 {
		pts = 10
	}
	totalPoints += pts
	if hasSub {
		earnedPoints += s.Score / 100.0 * pts
	}
}
pct := 0.0
if totalPoints > 0 {
	pct = earnedPoints / totalPoints * 100.0
}
\end{lstlisting}

Three behaviours in this loop are subtle and each is deliberate:

\begin{itemize}[leftmargin=1.4em]
  \item \textbf{Skipped questions stay in the denominator at zero.} Because the loop iterates the
  exam's problem list rather than the student's submissions, a problem the student never answered
  still adds its points to \texttt{totalPoints} while contributing nothing to
  \texttt{earnedPoints}. A student cannot inflate their percentage by leaving questions blank.
  \item \textbf{Pending-review answers are excluded from both numerator and denominator.} A
  written answer awaiting manual grading is \texttt{continue}d over entirely, so an ungraded essay
  does not depress the visible score to look like a failure before the teacher has read it. Once
  the teacher grades it and re-runs the aggregator, it re-enters both sums.
  \item \textbf{Points weighting happens here, at the problem level.} Each problem contributes in
  proportion to its \texttt{Points} (defaulting to 10 when unset or non-positive), so a 20-point
  hard problem and a 5-point easy one weigh proportionally. The submission's own \texttt{score} is
  a 0--100 percentage; scaling it by \texttt{pts / 100} converts it into earned points.
\end{itemize}

The computed percentage is written to \texttt{ExamAttempt.score}. Because the whole score is
recomputed from the children every time, the function is idempotent and safe to call any number
of times---there is no running total to double-count, and manual grading simply re-runs it.

\subsection{The single, deduplicated notification}

The final stage fires exactly one ``Exam Graded'' notification, and only under the right
conditions. The aggregator counts any submissions still in \texttt{pending\_review}; if written
answers remain ungraded, it updates the score but withholds the notification. Only when no
\texttt{pending\_review} submissions remain, the attempt's \texttt{status} is \texttt{submitted},
and the attempt has not already been notified does it fire---using the boolean column
\texttt{exam\_attempts.graded\_notified} as an idempotency flag, set true before the notification
is created so it can never be sent twice.

\begin{lstlisting}[language=Go,caption={Once-only ``Exam Graded'' notification, deduplicated by \texttt{graded\_notified} (\texttt{grader.go}).},label={lst:notify}]
if attempt.Status != "submitted" || attempt.GradedNotified {
	return
}
database.DB.Model(&attempt).Update("graded_notified", true)
notification.Create(attempt.UserID, "result",
	"Exam Graded",
	fmt.Sprintf("Your submission for \"%s\" has been graded. Score: %.0f%%", attempt.Exam.Title, pct),
	fmt.Sprintf("/dashboard/exam/%d/review", attempt.ExamID),
)
\end{lstlisting}

The result is a clean invariant: no matter how many questions an exam contains or in what order
their goroutines finish, the student receives one submission notification when they submit and one
grading notification when the last answer---automatic or manual---is graded.

\section{Sample versus Hidden Test Cases}\label{sec:grading-samplehidden}

The distinction between sample and hidden test cases runs through the whole pipeline and is what
prevents students from reverse-engineering the answer key. Each \texttt{TestCase} carries an
\texttt{is\_sample} boolean, and the three code paths treat it differently:

\begin{itemize}[leftmargin=1.4em]
  \item \textbf{Sample tests are visible.} The ``Run sample tests'' path
  (\texttt{RunSolution}) filters \texttt{is\_sample = true} and returns input, expected output,
  and actual output, giving students pre-submission feedback.
  \item \textbf{Hidden tests are opaque.} The real grader (\texttt{Grade}) applies no
  \texttt{is\_sample} filter and runs every test, so hidden tests count toward the score but are
  never surfaced during the exam.
  \item \textbf{Results reveal only samples.} On the results screen,
  \texttt{GetSubmission} attaches \texttt{input} and \texttt{expected\_output} to a result only
  when its test case is a sample; for hidden tests the student sees pass/fail and timing but never
  the input or the expected output.
\end{itemize}

Test-case management is teacher-owned and side-effecting. \texttt{AddTestCase},
\texttt{UpdateTestCase}, and \texttt{DeleteTestCase} in \texttt{internal/testcase/handler.go}
each verify ownership---the teacher must own the bank problem, or own the exam the problem belongs
to---before mutating a test. Critically, any change that could alter grading (adding a test,
deleting a test, or editing a test's input or expected output on a problem attached to an exam)
triggers \texttt{examsvc.ResetExamAttempts(*problem.ExamID)}, so that prior attempts graded
against the old test set are invalidated rather than left silently inconsistent. This connects the
grading engine to the destructive-edit reset mechanism described in \cref{ch:ch04}.

\section{The Honest Gap: Fire-and-Forget Durability}\label{sec:grading-queue}

The most significant limitation of the grading engine is the one already foreshadowed: grading is
dispatched as a bare goroutine (\texttt{go judge0.Grade(sub.ID)}) with no durable queue, no
supervisor, and no retry. The failure mode is precise and worth stating exactly. \texttt{Grade}
flips the submission to \texttt{running} before contacting Judge0; if the backend process crashes
or is redeployed while a grade is in flight, the goroutine dies with the process and the row is
left permanently in \texttt{running}. Because \texttt{maybeFinalizeAttempt} refuses to proceed
while any sibling is \texttt{pending} or \texttt{running}, the whole attempt never finalises and
the student never receives their grade. There is also no concurrency cap: a large exam could spawn
hundreds of simultaneous Judge0 requests and exhaust the public endpoint's rate limit. This gap is
documented openly in the project's README.

The mitigating context is scale. At the concurrency of a single classroom---tens of students, not
thousands---the model works, and the simplicity of returning the HTTP response instantly while
grading proceeds in the background is a real benefit. The gap is a considered trade-off, not an
accident.

The fix is well understood and follows directly from a property the current design already has:
\textbf{grading is idempotent}. \texttt{Grade} recomputes every result from the submission's stored
\texttt{code} and the problem's test cases and overwrites the previous results, so re-running it is
always safe. The durable version replaces the bare goroutine with the enqueue of a job into a
persistent work queue---a database-backed jobs table, or a broker such as Redis---consumed by a
bounded pool of workers. On start-up, a recovery sweep re-enqueues anything stuck in
\texttt{running}, and because re-grading is idempotent, the replay is safe. The bounded pool
simultaneously caps concurrency against Judge0. None of this changes the grading logic in this
chapter; it changes only how that logic is scheduled and made durable, which is why it sits on the
roadmap rather than blocking the current classroom-scale deployment.

\section{Summary}\label{sec:grading-summary}

This chapter traced a code submission from the Monaco editor to a graded, notified exam attempt.
The design rests on a clean separation of concerns: execution is delegated to
Judge0~\cite{judge0,dosilovic2018judge0} and its \texttt{isolate} sandbox~\cite{isolate}, so
untrusted code never touches the application host, while the \emph{judging}---output
normalisation, pass/fail determination, proportional scoring, and points-weighted
aggregation---is owned entirely by APEX. The nine-line \texttt{normalizeOutput} routine is the
fairness keystone, applied symmetrically to actual and expected output so that grading measures
content rather than whitespace~\cite{douce2005autoassess}. Around the coding grader, MCQ answers
are auto-judged by set equality and written answers enter a manual-grading queue, and all three
formats aggregate through \texttt{maybeFinalizeAttempt} into a single attempt score and a single,
deduplicated notification. The one material weakness---fire-and-forget dispatch without a durable
queue---is disclosed plainly, along with an idempotency-based fix that the existing design already
makes straightforward. Together these pieces deliver the project's central promise: three answer
formats, one grading pipeline, one score.
